Se ejecutaron dos escenarios en Docker Desktop para Linux. El primero partió de una instantánea JSON del catálogo obtenida con el conector de Supabase: una ciudad, trece direcciones, una zona y trece plazas. Las tablas de dispositivos e histórico estaban vacías. Se creó una copia local mediante pg\_dump, se simularon las pérdidas y la API devolvió cero plazas. Tras pg\_restore, la API volvió a devolver trece plazas y se verificaron las seis tablas y sus seis claves foráneas.

El segundo escenario utilizó tres plazas ficticias, un dispositivo y un registro histórico. La recuperación conservó también esas relaciones y los estados de la API. Se rechazaron una copia corrupta, rutas fuera de backups, una segunda operación simultánea y un destino sin la marca local. Un fallo SQL provocado antes de COMMIT revirtió la restauración. Dos pruebas de retención confirmaron que solo se eliminan copias válidas antiguas y que una configuración inválida no elimina archivos.

La primera prueba tardó aproximadamente 3,1 segundos entre la confirmación de pérdida y la verificación final. Este resultado corresponde al pequeño conjunto local y no constituye una garantía para producción. Los eventos reales seleccionados se guardan en \texttt{evidencias.jsonl}; las copias, los logs completos y la captura del mapa permanecen en el portátil, excluidos de Git.

No se ha probado el pg\_dump remoto de Supabase porque falta su conexión PostgreSQL. La instantánea JSON no conserva todo el DDL original. La copia remota, su revisión de RLS, la tarea cron y la copia cifrada fuera del equipo quedan pendientes. La memoria se entrega como fuente LaTeX sin compilación en este entorno. No se modificaron Supabase ni el despliegue activo de Vercel.
